\section{Methodology}

\subsection{Control Objective}

The governance objective combines throughput, latency, decentralization proxy, and target tracking:
\[
\begin{aligned}
r ={}& w_1 \cdot \text{throughput} + w_2 \cdot \text{latency}(d) \\
&+ w_3 \cdot \text{decentralization} - \lambda |d - d_t|.
\end{aligned}
\]
Action outputs are bounded and then validated by a safety shield before execution.

\subsection{Policy Family}

We evaluate random, heuristic, PPO-family, and robust gated controllers. PPO uses clipped updates \cite{schulman2017ppo}. The robust branch adds risk shaping:
\[
\tilde r_t = r_t - \lambda_{\mathrm{risk}}\,\mathrm{Var}(r_{t-k:t}).
\]

\subsection{Shield Logic (Executable Rule)}

The shield is deterministic and state-aware. At each step, it enforces bounded actions and bounded action-rate change before execution:
\[
a^{\mathrm{exec}}_t=\mathrm{clip}(a_t, a_{\min}, a_{\max}),
\]
\[
\Delta a_t=\mathrm{clip}(a^{\mathrm{exec}}_t-a_{t-1},-\delta_{\max},\delta_{\max}).
\]
If a post-check still violates a hard safety predicate (for example, projected delay or partition threshold), the shield routes to deterministic fallback $a_t^{\mathrm{heur}}$ and re-validates. In implementation terms:
\begin{itemize}
    \item If \texttt{hard\_safe(state, action)} is false, clamp action by bounds and max delta.
    \item Re-check safety; if still false, switch to \texttt{heuristic\_action(state)}.
    \item Return the validated action.
\end{itemize}

For explicit auditability, shield execution is summarized below:
\begin{enumerate}
    \item Input $(s_t, a_t, a_{t-1})$.
    \item Set $a\leftarrow\mathrm{clip}(a_t,a_{\min},a_{\max})$.
    \item Set $a\leftarrow a_{t-1}+\mathrm{clip}(a-a_{t-1},-\delta_{\max},\delta_{\max})$.
    \item If $\mathrm{hard\_safe}(s_t,a)$, return $a$.
    \item Set $a_h\leftarrow\mathrm{heuristic\_action}(s_t)$.
    \item Set $a_h\leftarrow\mathrm{clip}(a_h,a_{\min},a_{\max})$.
    \item Set $a_h\leftarrow a_{t-1}+\mathrm{clip}(a_h-a_{t-1},-\delta_{\max},\delta_{\max})$.
    \item If $\mathrm{hard\_safe}(s_t,a_h)$, return $a_h$.
    \item Otherwise, return $\mathrm{emergency\_safe\_action}(s_t)$.
\end{enumerate}

\subsection{Quantum-Inspired Uncertainty Gate}

The policy front-end is a 6-qubit variational circuit with angle embedding, entangling layers, data re-uploading, and Pauli-$Z$ measurements for mean/scale signals. Uncertainty is estimated by parameter perturbation:
\[
\begin{aligned}
u_q(x) &= \operatorname{Var}\left(\{\mu_{\theta+\epsilon_k}(x)\}_{k=1}^{K}\right), \\
\epsilon_k &\sim \mathcal{N}(0,0.01I),\quad K=10.
\end{aligned}
\]
The routing rule is
\[
a_t=
\begin{cases}
a_t^{\mathrm{heur}}, & u_q(x_t)>\tau_q,\\
a_t^{\mathrm{policy}}, & u_q(x_t)\le\tau_q.
\end{cases}
\]

\subsection{Engineering Trade-offs in Uncertainty Design}

Latency mattered when we picked the uncertainty block, not only whether the score looked pretty offline. Full ensembles diversify but you pay with memory and serving: every step needs several forward passes through separate copies. Parameter perturbation keeps one small network and jitters weights at inference; that is easier to host next to a real controller.

Head-to-head, the perturbation readout hit higher precision at the same trigger budget than MC-dropout and the ensemble baselines we tried, without giving up much recall. We therefore tuned for ``good enough routing under a fixed trigger rate'' instead of chasing a single scalar uncertainty metric.

\begin{figure}[H]
\centering
\includegraphics[width=0.88\linewidth]{figures/robust_policy_flow.pdf}
\caption{Quantum-inspired uncertainty-gated fallback and shielded action execution.}
\label{fig:robust-flow}
\end{figure}

\subsection{Why the Compact Estimator Compresses Well Here}

We simulate the six-qubit VQC-style block classically; the win we cared about was parameter count, not quantum throughput. Governance states in our build mix a handful of heavy drivers---delay, congestion, attack pressure---with ugly cross-terms. Entanglement buys multiplicative mixing cheaply, and re-uploading lets the same wires react nonlinearly across time. That recipe tracks the control surface we needed with far fewer weights than a fat MLP. In practice the compression only looked good once gating and shielding were on, because otherwise the head saw regimes it was never trained to stabilize.

\subsection{Adversarial and Ablation Design}

The scaled benchmark injects burst demand, delay noise, attack pressure, and regime shifts. We evaluate seven configurations:
\begin{itemize}
    \item full system (quantum gate + shield + robust policy),
    \item no quantum gate,
    \item no shield,
    \item no fallback,
    \item classical PPO,
    \item Heuristic controller,
    \item random.
\end{itemize}
The stress-aware reward is
\[
\begin{aligned}
r_t^{\text{adv}} ={}& \alpha\,\text{throughput}_t - \beta\,\text{delay}_t - \gamma\,\text{attack}_t \\
&- \eta\,\text{fork}_t - \lambda |d_t-d^\star|.
\end{aligned}
\]

\subsection{Evaluation}

We track the obvious scalars---reward, violations, target MAE, how often fallback fires, and a latency proxy. Whenever we compare PPO-family runs we pair that with bootstrap intervals, permutation tests, and Cohen's $d$ so we do not overfit a table to one lucky seed \cite{efron1994introduction,good2005permutation}.